%! AP Statistics — Unit 1, Lesson 1.8 — Warm-Up ANSWER KEY
\documentclass[10pt]{article}
\usepackage{apstats-article}
\usepackage{apstats-key}

%=============================================================================
\begin{document}

\pageheader{Unit 1, Lesson 1.8}{Warm-Up --- Answer Key}
\begin{tabularx}{\linewidth}{X X X}
Name:\;\ans{Key} & Date:\; & Period:\;
\end{tabularx}

\vspace{0.16in}

\noindent In Lesson 1.7 you learned to calculate the five-number summary and the
$1.5\times\text{IQR}$ fence rule. Today you will see how those five numbers
create a picture of a distribution. Use the boxplot below, which displays the
\textbf{hours spent on homework per week} for a sample of 20 high school students.

\vspace{0.18in}

\begin{center}
\begin{tikzpicture}[scale=1.0]
  \draw[thick] (0.0,0) -- (7.5,0);
  \foreach \v/\lbl in {0/0, 0.52/1, 1.04/2, 1.56/3, 2.08/4, 2.60/5,
                        3.12/6, 3.64/7, 4.16/8, 4.68/9, 5.20/10,
                        5.72/11, 6.24/12} {
    \draw (\v,-0.12) -- (\v,0.12);
    \node[below, font=\small] at (\v,-0.12) {\lbl};
  }
  \node[below, font=\small\itshape] at (3.12,-0.56) {Homework hours per week};
  \draw[thick, navy] (0.52, 0.38) -- (1.56, 0.38);
  \draw[thick, navy] (1.56, 0.10) rectangle (4.68, 0.66);
  \draw[thick, navy] (3.12, 0.10) -- (3.12, 0.66);
  \draw[thick, navy] (4.68, 0.38) -- (6.24, 0.38);
  \draw[thick, navy] (0.52, 0.22) -- (0.52, 0.54);
  \draw[thick, navy] (6.24, 0.22) -- (6.24, 0.54);
\end{tikzpicture}
\end{center}

\vspace{0.22in}

\begin{enumerate}[leftmargin=1.5em, itemsep=8pt]

  \item What are the minimum and maximum values shown in the boxplot?
        What is the \textbf{range}?\\[4pt]
        \ans{Minimum $= 1$ hour; Maximum $= 12$ hours.
        Range $= 12 - 1 = \mathbf{11}$ hours.}

  \item Identify $Q_1$, the \textbf{median}, and $Q_3$ from the boxplot.
        Then calculate the \textbf{IQR}.\\[4pt]
        \ans{$Q_1 = 3$, $M = 6$, $Q_3 = 9$.
        IQR $= 9 - 3 = \mathbf{6}$ hours.}

  \item Calculate the \textbf{upper fence}. Would a student reporting 14 hours
        be flagged as an outlier? Show your work.\\[4pt]
        \ans{Upper fence $= Q_3 + 1.5(\text{IQR}) = 9 + 1.5(6) = 9 + 9 = \mathbf{18}$ hours.
        Since $14 < 18$, the value 14 is \textbf{not} an outlier.}

  \item Looking at the boxplot, is this distribution \textit{symmetric},
        \textit{skewed left}, or \textit{skewed right}? Explain using the
        positions of the median and the whiskers.\\[4pt]
        \ans{The distribution is approximately \textbf{symmetric}. The median (6) is
        centered within the box, and the two whiskers are equal in length (each spans
        2 hours). There is no strong pull toward either tail.}

\end{enumerate}

\vspace{0.14in}

\begin{tcolorbox}[colback=greenbg, colframe=greenacc, boxrule=0.8pt, arc=2mm,
    left=4mm, right=4mm, top=3mm, bottom=3mm]
\small
\begin{itemize}[leftmargin=1.3em, itemsep=5pt]
  \item Today you will learn how boxplots are \textit{built} from the five-number
        summary --- one step at a time.
  \item You will also see how a \textbf{modified boxplot} handles outliers
        differently from a standard boxplot.
  \item Key question to keep in mind: What can a boxplot show that a histogram
        cannot, and vice versa?
\end{itemize}
\end{tcolorbox}

\begin{teachernote}
\textbf{Grading notes (4 pts total):}
\begin{itemize}[itemsep=2pt]
  \item Q1: 1 pt — min=1, max=12, range=11 all required.
  \item Q2: 1 pt — all three values correct ($Q_1=3$, $M=6$, $Q_3=9$, IQR=6).
  \item Q3: 1 pt — correct fence calculation AND correct outlier conclusion (14 is NOT an outlier).
  \item Q4: 1 pt — names symmetric AND provides two pieces of evidence (median placement + equal whiskers).
\end{itemize}
\end{teachernote}

\end{document}
